\section{案例研究：灵巧操作中的 Model-Based RL}

讲者介绍了一个具体案例，使用 model-based RL 让机器人灵巧手进行物体操作。

\begin{figure}[H]
\centering
\includegraphics[width=0.9\textwidth]{slides-images/slide-35.jpg}
\caption{灵巧操作的模拟消融实验结果\protect\footnotemark}
\end{figure}
\footnotetext{Slides 第 35 页。}

关键实验发现：
\begin{itemize}
    \item 需要足够大的模型架构（2x250 或 2x500 隐藏单元）
    \item 至少需要 3 个集成成员来获得可靠的不确定性估计
    \item 规划视野存在权衡：太短看不到足够远的未来，太长模型误差累积
    \item 改进的 CEM 控制器（如 Filtering + Reward Weighting）比 Random Shooting 好得多
\end{itemize}

\subsection{何时使用 Model-Based RL}

\begin{importantbox}{Model-Based RL 的适用场景}
\begin{itemize}
    \item 数据收集代价高昂（如真实机器人）
    \item 状态空间相对低维且可预测
    \item 需要快速适应新任务（模型可迁移）
    \item 动力学具有一定的可学习结构
\end{itemize}
不适合的场景：高维观测（如图像）且动力学高度随机的环境。
\end{importantbox}

\subsection{本章小结}

灵巧操作案例展示了 model-based RL 在数据效率上的优势，以及集成、规划视野等超参数对性能的关键影响。

